% ===========================================================================
%  LiteraLevel — Presentación (Módulo 5, 10 puntos)
%  Diplomado en Ciencia de Datos, FES Acatlán (UNAM), Generación 34
%
%  Mismos colores, misma tipografia y mismos escudos que documento.tex y
%  guion_video.tex: los tres entregables tienen que parecer del mismo trabajo.
%  Las figuras se dibujan en TikZ, no se pegan como imagen.
%
%    tectonic presentacion.tex
% ===========================================================================

\documentclass[aspectratio=169,11pt]{beamer}

\usepackage{iftex}
\ifPDFTeX
  \usepackage[T1]{fontenc}
  \usepackage[utf8]{inputenc}
  \usepackage{lmodern}
\else
  \usepackage{fontspec}
\fi

\usepackage[spanish,es-noquoting,es-nodecimaldot]{babel}
\usepackage{graphicx}
\usepackage{booktabs}
\usepackage{array}
\usepackage{tikz}
\usepackage{pgfplots}
\pgfplotsset{compat=1.18}
% Coma decimal en las graficas, igual que en el documento.
\pgfkeys{/pgf/number format/.cd, use comma, 1000 sep={\,}}
\usetikzlibrary{positioning}

% --- La misma paleta que el documento -------------------------------------
\definecolor{llazul}{HTML}{1F4E79}
\definecolor{llcian}{HTML}{2E86AB}
\definecolor{llverde}{HTML}{2A9D8F}
\definecolor{llambar}{HTML}{E9A13B}
\definecolor{llrojo}{HTML}{C1495B}
\definecolor{llgris}{HTML}{5A5A5A}
\definecolor{llgrisclaro}{HTML}{E8ECF0}
\definecolor{llfondo}{HTML}{F7F9FB}

\usetheme{default}
\usecolortheme{default}
\setbeamertemplate{navigation symbols}{}
\setbeamercolor{frametitle}{fg=llazul}
\setbeamercolor{title}{fg=llazul}
\setbeamercolor{structure}{fg=llcian}
\setbeamercolor{normal text}{fg=black}
\setbeamerfont{frametitle}{series=\bfseries}
\setbeamertemplate{itemize item}{\color{llcian}\textbullet}
\setbeamertemplate{itemize subitem}{\color{llgris}--}

\setbeamertemplate{footline}{%
  \hfill\raisebox{2mm}{\color{llgris}\scriptsize
    LiteraLevel · Módulo 5 \quad \insertframenumber\hspace{6mm}}}

% --- Una cifra grande y sola, que es la regla del guion --------------------
\newcommand{\cifra}[2]{%
  \begin{minipage}{\linewidth}\centering
    {\fontsize{42}{46}\selectfont\bfseries\color{llazul}#1}\\[4pt]
    {\small\color{llgris}#2}
  \end{minipage}}

\newcommand{\cod}[1]{{\ttfamily\small #1}}

\title{LiteraLevel}
\subtitle{¿Está este libro en inglés a tu nivel de lectura?}
\author{Miguel Reyes Carrillo}
\date{Septiembre de 2026}

\begin{document}

% ===========================================================================
\begin{frame}[plain]
\vspace{2mm}
\begin{center}
\begin{minipage}{0.16\textwidth}\raggedright
  \includegraphics[height=1.6cm]{logos/logo_unam.png}
\end{minipage}
\hfill
\begin{minipage}{0.62\textwidth}\centering
  {\small\bfseries Universidad Nacional Autónoma de México}\\[2pt]
  {\scriptsize FACULTAD DE ESTUDIOS SUPERIORES ACATLÁN}
\end{minipage}
\hfill
\begin{minipage}{0.16\textwidth}\raggedleft
  \includegraphics[height=1.6cm]{logos/logo_fes_acatlan.png}
\end{minipage}
\end{center}

\vspace{6mm}
\begin{center}
  {\Huge\bfseries\color{llazul} LiteraLevel}\\[6pt]
  {\large ¿Está este libro en inglés a tu nivel de lectura?}\\[10mm]
  {\normalsize Miguel Reyes Carrillo}\\[3pt]
  {\small\color{llgris} Proyecto integrador --- Módulo 5 \quad·\quad Generación 34}\\[2pt]
  {\small\color{llgris} Mtro. Fernando Barranco Rodríguez}
\end{center}
\end{frame}

% ===========================================================================
\begin{frame}{Antes de comenzar}
\vspace{4mm}
\begin{center}
{\large Un hispanohablante elige un libro en inglés que suena bien,\\[2pt]
lo abre\ldots{} y a las tres páginas lo cierra.}\\[8mm]

{\Large\bfseries\color{llazul} Quizá no era falta de capacidad.}\\[3pt]
{\Large\bfseries\color{llazul} Quizá necesitaba otra lectura.}\\[10mm]

{\color{llgris} Nadie le dijo antes de abrirlo si estaba a su altura.}
\end{center}
\end{frame}

% ===========================================================================
\begin{frame}{Peldaño}
\vspace{3mm}
\begin{center}
{\large La compañía —\textbf{ficticia}, como pide el módulo— que hace esto.}\\[2pt]
{\small\color{llgris} El modelo, la aplicación y todas las cifras que siguen no lo son.}\\[7mm]
\end{center}

\begin{columns}[T]
\begin{column}{0.48\textwidth}
\textbf{\color{llazul} El nicho}\\[2pt]
{\small Lectura extensiva en inglés para hispanohablantes \textbf{entre A2 y B2}:
ya no caben en el material graduado y todavía no aguantan una novela sola.}\\[5mm]

\textbf{\color{llazul} El nombre}\\[2pt]
{\small Se aprende con lo que está \textbf{un escalón} por encima de lo que uno
domina. Encontrar ese peldaño es el problema.}
\end{column}
\begin{column}{0.48\textwidth}
\textbf{\color{llazul} Quién paga qué}\\[2pt]
{\small Medir un libro, leerlo, buscar y el diccionario: \textbf{gratis}, corren
en local sobre datos abiertos.\\[3pt]
Explicar el contexto, traducir, resumir y corregir la pronunciación: usan
OpenAI, con permiso y cuota.\\[3pt]
Sin clave, \textbf{la lectura y el diccionario local siguen disponibles}.}
\end{column}
\end{columns}
\end{frame}

% ===========================================================================
\begin{frame}{Los datos, en cifras}
\vspace{6mm}
\begin{columns}[T]
\begin{column}{0.33\textwidth}\centering
  \cifra{4.724}{extractos con dificultad\\medida por lectores (CLEAR)}
\end{column}
\begin{column}{0.33\textwidth}\centering
  \cifra{654}{documentos anclados\\a niveles del MCER}
\end{column}
\begin{column}{0.33\textwidth}\centering
  \cifra{61.642}{libros objetivo en inglés\\Gutenberg, 7 sep. 2026}
\end{column}
\end{columns}

\vspace{10mm}
\begin{center}\small\color{llgris}
Ninguno viene de Kaggle: catálogo oficial, repositorio público y descarga
directa de los PDF.
\end{center}
\end{frame}

% ===========================================================================
\begin{frame}{Objetivos}
\vspace{4mm}
\begin{itemize}\setlength\itemsep{5mm}
  \item[{\color{llcian}\bfseries 01}] \textbf{Situar al lector} en una escala
        continua, con un cuestionario breve o midiendo un texto que ya lee cómodo.
  \item[{\color{llcian}\bfseries 02}] \textbf{Medir un libro entero}, capítulo
        a capítulo, y no una muestra de sus primeras páginas.
  \item[{\color{llcian}\bfseries 03}] \textbf{Recomendar qué leer hoy}, con la
        dificultad puesta al lado del nivel del lector.
  \item[{\color{llcian}\bfseries 04}] \textbf{Acompañar la lectura}: diccionario,
        explicación contextual y práctica de pronunciación.
\end{itemize}
\end{frame}

% ===========================================================================
\begin{frame}{Hallazgo 1 · El cuello de botella era la representación}
\vspace{2mm}
{\small Con las mismas 950 muestras reservadas, cambiando solo cómo se
representa el texto:}\\[1mm]
{\scriptsize\color{llgris} Esas 950 muestras eligieron además la época del
transformador: son validación de selección, no un test independiente. El error
ordinal es sobre sextiles de CLEAR, no sobre niveles MCER.}
\vspace{2mm}
\begin{center}
\begin{tikzpicture}
\begin{axis}[
  width=0.86\linewidth, height=4.4cm, xbar, bar width=9pt,
  xmin=0, xmax=1.85, xlabel={Error ordinal medio (menos es mejor)},
  symbolic y coords={trans,emb,conteos},
  ytick={conteos,emb,trans},
  yticklabels={3 variables clásicas,\emph{Embedding} de 384 d,Transformador},
  yticklabel style={font=\scriptsize}, tick style={draw=none},
  xmajorgrids=true, grid style={llgrisclaro},
  axis line style={draw=llgrisclaro}, enlarge y limits=0.3,
  nodes near coords={\pgfmathprintnumber[fixed,precision=2]{\pgfplotspointmeta}},
  nodes near coords style={font=\scriptsize, color=llgris},
  nodes near coords align={horizontal},
]
\addplot[fill=llgris!25, draw=llgris!60] coordinates {(1.39,conteos)};
\addplot[fill=llcian!35, draw=llcian!80] coordinates {(0.93,emb)};
\addplot[fill=llazul!45, draw=llazul!90] coordinates {(0.67,trans)};
\end{axis}
\end{tikzpicture}
\end{center}

\vspace{-1mm}
{\footnotesize Con las mismas variables, un XGBoost de 300 árboles queda por
detrás de una regresión logística. Cambiar la representación sí mejora el
resultado. Las tres barras proceden de \textbf{la misma tanda}.}
\end{frame}

% ===========================================================================
\begin{frame}{Hallazgo 2 · Contar palabras llega casi igual de lejos}
\vspace{5mm}
\begin{columns}[T]
\begin{column}{0.5\textwidth}\centering
  \cifra{0,860}{el modelo, contra el nivel\\oficial de 47 textos}
\end{column}
\begin{column}{0.5\textwidth}\centering
  \cifra{0,807}{\textbf{contar las palabras},\\sin modelo alguno}
\end{column}
\end{columns}

\vspace{9mm}
{\small En un corpus de exámenes, el nivel y la extensión de la tarea van
juntos, así que un contador de palabras aprovecha un atajo ligado al diseño de
esos exámenes, cuya utilidad fuera de ese corpus no está demostrada. \textbf{Lo que sostiene al modelo es lo que aporta por
encima de la longitud}: descontándola, conserva \textbf{0,668}.}

\vspace{3mm}
{\footnotesize\color{llgris} Se enseña porque es la comparación más
desfavorable del proyecto, y ocultarla sería el peor error posible.}
\end{frame}

% ===========================================================================
\begin{frame}{El modelo desplegado}
\vspace{4mm}
\begin{center}
{\large\bfseries\color{llazul} distilroberta-base} \quad·\quad
82 M de parámetros \quad·\quad regresión, no clasificación\\[6mm]

\begin{tabular}{@{}l r r r@{}}
\toprule
\textbf{Conjunto} & \textbf{n} & \textbf{Dentro de $\pm 1$ nivel} & \textbf{Spearman} \\
\midrule
Validación interna & 144 & 97,2\,\% & 0,908 \\
British Council, desarrollo & 38 & 100\,\% & 0,829 \\
\bottomrule
\end{tabular}\\[3mm]
{\scriptsize\color{llgris} Métricas internas recalculadas el 12 de septiembre con el
guion de entrenamiento y datos públicos: idénticas a cuatro decimales.}
\\[3mm]

{\small Sobre textos externos con nivel oficial, las medias por nivel salen
\textbf{estrictamente crecientes}:}\\[3mm]
{\small A1 \textbf{0,504} \; · \; A2 \textbf{0,556} \; · \; B1 \textbf{0,928}
\; · \; B2 \textbf{1,241} \; · \; C1 \textbf{1,495}}
\end{center}

\vspace{3mm}
{\footnotesize\color{llgris} El orden es el resultado, no el porcentaje. El
escalón A1--A2 descansa en 3 textos y está dentro del ruido; los otros tres,
no.}
\end{frame}

% ===========================================================================
\begin{frame}{La aplicación entregada}
\vspace{1mm}
{\small Angular es la interfaz del lector; la API conecta con el motor de Python.}
\vspace{2mm}

\begin{center}
\begin{tikzpicture}[font=\scriptsize, node distance=5mm]
\tikzset{caja/.style={draw=llcian!70, fill=llfondo, rounded corners=2pt,
                      text width=3.4cm, align=center, inner sep=4pt,
                      minimum height=1.5cm}}
\node[caja] (f) {\textbf{Front · Angular 22}\\ Interfaz entregada\\
  Space estático\\ \emph{Leer · Encontrar · Pronunciar · Tu nivel}};
\node[caja, right=of f] (a) {\textbf{API · FastAPI}\\ Rutas bajo \cod{/api}\\
  reutiliza los módulos\\ y las puertas de pago\\ del motor};
\node[caja, right=of a] (m) {\textbf{Motor}\\ modelo v4 · diccionario\\
  Space de Gradio,\\ \cod{cpu-basic}; Gradio\\ en la raíz, de respaldo};
\draw[->, llcian!80, thick] (f) -- (a);
\draw[->, llcian!80, thick] (a) -- (m);
\end{tikzpicture}
\end{center}

\vspace{2mm}
\begin{columns}[T]
\begin{column}{0.32\textwidth}
\textbf{\color{llazul} Con una mano}\\[2pt]
{\small Barra de lectura al alcance del pulgar y acceso al micrófono desde
la vista de pronunciación.}
\end{column}
\begin{column}{0.32\textwidth}
\textbf{\color{llazul} Voz e IPA}\\[2pt]
{\small Compara lo dicho con la frase objetivo; ofrece IPA de diccionario,
audio y orientación para practicar el sonido.}
\end{column}
\begin{column}{0.32\textwidth}
\textbf{\color{llazul} Apoyo articulatorio}\\[2pt]
{\small Modelo Blender de 38 piezas en la ficha IPA y la corrección.
El 3D se carga al abrirlo, con respaldo 2D. Apoyo didáctico, no clínico.}
\end{column}
\end{columns}

\vspace{3mm}
{\footnotesize\color{llgris} Las portadas FLUX siguen pendientes;
la demostración no depende de imágenes aún no generadas.}
\end{frame}

% ===========================================================================
\begin{frame}{Lo que hace por el lector}
\vspace{1mm}
{\small Lo añadido en los dos últimos días, ya publicado en el \emph{front}.}
\vspace{3mm}

\begin{columns}[T]
\begin{column}{0.32\textwidth}
\textbf{\color{llazul} Ver la boca mientras suena}\\[2pt]
{\small El corte de perfil y la boca 3D se mueven al compás del audio del sonido
o de la palabra, también en modo lento. La lengua se distingue del paladar y un
aro marca dónde toca.}
\end{column}
\begin{column}{0.32\textwidth}
\textbf{\color{llazul} Leer con naturalidad}\\[2pt]
{\small Mide dónde se paró el lector y qué tramos van lentos para su nivel.
Marca énfasis, enlaces (\emph{with us}) y reducciones
(\emph{don't you} \textrightarrow{} \emph{doncha}), y lo hace escuchar dicho de
corrido.}
\end{column}
\begin{column}{0.32\textwidth}
\textbf{\color{llazul} Más lectura a su nivel}\\[2pt]
{\small \textbf{15.357} textos de licencia libre medidos con el v4
---StoryWeaver, Let's Read, African Storybook y VOA--- junto a los
\textbf{53.179} libros de Gutenberg.}
\end{column}
\end{columns}

\vspace{4mm}
{\footnotesize\color{llgris} Las reducciones solo salen de una lista cerrada y
los enlaces se comprueban con la pronunciación de diccionario; lo que escribe la IA
va marcado.}
\end{frame}

% ===========================================================================
\begin{frame}{Ejemplo de uso y aplicación}
\vspace{1mm}
{\small \textbf{Caso ilustrativo del documento}: perfil \textbf{1,20 sobre 2},
banda intermedia. No es una sesión registrada ni un resultado fijo del cuestionario.}
\vspace{2mm}

\begin{center}
\begin{tikzpicture}[font=\scriptsize, node distance=2mm]
\tikzset{p/.style={draw=llcian!70, fill=llfondo, rounded corners=2pt,
                   text width=2.25cm, align=center, inner sep=3.5pt,
                   minimum height=1.1cm}}
\node[p] (a) {\textbf{1. Se sitúa}\\ cuestionario\\ $\rightarrow$ 1,20};
\node[p, right=of a] (b) {\textbf{2. Busca}\\ escribe mal\\ «dostoyevski»};
\node[p, right=of b] (c) {\textbf{3. Abre}\\ \emph{Crime and}\\ \emph{Punishment}};
\node[p, right=of c] (d) {\textbf{4. Compara}\\ curva guardada\\ con su nivel};
\node[p, right=of d] (e) {\textbf{5. Lee}\\ diccionario\\ y voz alta};
\foreach \x/\y in {a/b,b/c,c/d,d/e}
  \draw[->, llcian!80, thick] (\x) -- (\y);
\end{tikzpicture}
\end{center}

\vspace{2mm}
\begin{center}\small
\begin{tabular}{@{}l l@{}}
\toprule
\textbf{Paso} & \textbf{Lo que produce} \\
\midrule
Se sitúa & 1,20 sobre 2, banda \emph{intermedia}, ancla B2 \\
Busca & búsqueda por significado; el nombre mal escrito llega igual \\
Abre & ficha sobre el texto: extensión, diálogo y repeticiones \\
Compara & ejemplo registrado: dificultad \textbf{1,70} frente a perfil \textbf{1,20} \\
Lee & diccionario primero; explicación de IA después, si está habilitada \\
\bottomrule
\end{tabular}
\end{center}

\vspace{1mm}
{\footnotesize\color{llgris} La medición histórica del libro no es una promesa
de latencia HTTP. En la demo se muestran los valores que devuelve la app.}
\end{frame}

% ===========================================================================
\begin{frame}{Contra quién compite}
\vspace{4mm}
\begin{center}
\begin{tikzpicture}
\begin{axis}[
  width=0.8\linewidth, height=3.9cm, xbar, bar width=11pt,
  xmin=0, xmax=1.02, xlabel={Pearson con el juicio humano, corpus CLEAR},
  symbolic y coords={ll,llm,form}, ytick={form,llm,ll},
  yticklabels={Fórmulas clásicas,LLM \emph{zero-shot},\textbf{Nuestro codificador}},
  yticklabel style={font=\scriptsize}, tick style={draw=none},
  xmajorgrids=true, grid style={llgrisclaro},
  axis line style={draw=llgrisclaro}, enlarge y limits=0.32,
  nodes near coords={\pgfmathprintnumber[fixed,precision=3]{\pgfplotspointmeta}},
  nodes near coords style={font=\scriptsize, color=llgris},
  nodes near coords align={horizontal},
]
\addplot[fill=llgris!25, draw=llgris!60] coordinates {(0.517,form)};
\addplot[fill=llgris!25, draw=llgris!60] coordinates {(0.725,llm)};
\addplot[fill=llcian!45, draw=llcian!90] coordinates {(0.89,ll)};
\end{axis}
\end{tikzpicture}
\end{center}

\vspace{-2mm}
{\footnotesize\textbf{No son comparaciones cara a cara}: las dos primeras
vienen de un estudio ajeno sobre el mismo corpus, con otra partición.
Justifican la decisión de afinar un codificador; no demuestran una victoria
medida.}

\vspace{3mm}
{\footnotesize\textbf{Y dónde se pierde:} Lexile tiene treinta años de
validación y una escala \textbf{que acepta un colegio}; la de aquí no la acepta
nadie.}
\end{frame}

% ===========================================================================
\begin{frame}{Estrategia}
\vspace{3mm}
{\small Acciones operativas y siguiente validación con lectores.}
\vspace{3mm}

\begin{center}\small
\begin{tabular}{@{}p{0.30\linewidth} p{0.30\linewidth} p{0.30\linewidth}@{}}
\toprule
\textbf{Acción} & \textbf{Recursos} & \textbf{Criterio de terminado} \\
\midrule
Servir interfaz y motor & Angular estático + API en Space de Gradio;
  IA con cuota &
  \textbf{Publicado}: comprobar navegación y respuesta JSON antes de la demo \\[3mm]
Ofrecer lecturas ajustadas & Catálogo medido por v4 y corpus graduados abiertos &
  \textbf{Disponible}: 53.179 libros de Gutenberg (585.111 capítulos) y 15.357 textos libres medidos \\[3mm]
Validar la utilidad & Pruebas con lectores y registro de tareas &
  \textbf{Pendiente}: medir comprensión y continuidad con usuarios \\
\bottomrule
\end{tabular}
\end{center}

\vspace{3mm}
{\footnotesize\color{llgris} La consulta local no llama a OpenAI. El servicio
remoto y la IA sí tienen costes operativos.}
\end{frame}

% ===========================================================================
\begin{frame}{Lo que esto no es}
\vspace{8mm}
\begin{center}
{\Large\bfseries\color{llrojo} No certificamos ningún nivel del MCER.}\\[6mm]

{\large La letra que se enseña es la etiqueta más cercana\\
en una escala continua propia.}\\[8mm]

{\normalsize\color{llgris} Sirve para elegir qué leer. No para acreditar nada.}
\end{center}
\end{frame}

% ===========================================================================
\begin{frame}[plain]
\vspace{16mm}
\begin{center}
{\Large\bfseries\color{llazul} Funcionando ahora mismo}\\[8mm]

{\large\ttfamily mikemaki1763-literalevel-front.static.hf.space}\\[6mm]

{\small\color{llgris} API y motor:
\ttfamily huggingface.co/spaces/MikeMaki1763/literalevel}\\[3mm]
{\small\color{llgris} Modelo público:
\ttfamily huggingface.co/MikeMaki1763/literalevel-v4}\\[10mm]

{\normalsize Elige una lectura, compara su dificultad con tu nivel\\[3pt]
y prueba si te ayuda a seguir leyendo.}
\end{center}
\end{frame}

\end{document}
